At every even horizon \(T=2n\), the bounded two-regime score and randomization process of Proposition~\(\text{prop:endogenous-qv-self-normalization-obstruction}\) has an explicit bounded augmentation with a nondegenerate observable-outcome law and a nonzero causal effect that satisfies every assumption explicitly listed in Saco's Theorem~5.14 and also its maintained consistency and no-interference condition. Namely, let \((U_t)_{t=1}^T\) be independent Rademacher variables, independent of the action randomizers, set \(A_t^{\mathrm S}:=\mathbf1\{A_t=1\}\) and \(X_t^{\mathrm S}:=(X_t,U_t)\), and set
\[
 Y_t^{\mathrm S}(1)=\frac32+U_t,
 \qquad
 Y_t^{\mathrm S}(0)=\frac12+U_t.
\]
At context value \((x,u)\), set the predictable nuisance regressions to \(\widehat m_{t-1,1}(x,u)=u+1/2\) and \(\widehat m_{t-1,0}(x,u)=u-1/2\). Then the causal target is \(\theta_0=E_P[Y_t^{\mathrm S}(1)-Y_t^{\mathrm S}(0)]=1\), while the observed outcome obeys \(\operatorname{Var}(Y_t^{\mathrm S}\mid\mathcal F_{t-1})=1+p_t(1-p_t)\ge1\), where \(p_t:=P(A_t=1\mid\mathcal F_{t-1})\). Saco's centered AIPW increment is nevertheless exactly \(\xi_t\). Conditionally on \(\mathcal F_{t-1}\),
\[
 P\!\left(\xi_t=\frac1{p_t}\,\middle|\,\mathcal F_{t-1}\right)=p_t,
 \qquad
 P\!\left(\xi_t=-\frac1{1-p_t}\,\middle|\,\mathcal F_{t-1}\right)=1-p_t,
\]
and hence \(E(\xi_t\mid\mathcal F_{t-1})=0\) and \(\operatorname{Var}(\xi_t\mid\mathcal F_{t-1})=1/p_t+1/(1-p_t)\in\{4,16/3\}\). Saco's realized and predictable quadratic variations are therefore \(Q_T=[S]_T\) and \(V_T^2=\langle S\rangle_T\), with \(V_T^2\ge4T\) almost surely. Nevertheless, along \(T=2n\),
\[
 \frac{S_T}{\sqrt{Q_T}}\Longrightarrow
 L:=\mathbf1\{Z>0\}\frac{2Z+4W/\sqrt3}{\sqrt{28/3}}
 +\mathbf1\{Z\le0\}\frac{2Z+2W}{\sqrt8},
\]
where \(Z,W\) are independent \(N(0,1)\) variables and
\[
 E[L]=\frac{2}{\sqrt{2\pi}}
 \left(\sqrt{\frac3{28}}-\frac1{\sqrt8}\right)\ne0.
\]
The feasible sample-variance Studentized AIPW statistic has the same limit \(L\). Thus the variance-growth-only normal-limit conclusions asserted in Saco's Theorem~5.14 are false as stated, even for bounded nondegenerate outcomes and a nonzero causal effect.